%%%%%%%%%%%%%%%%%%%%%%%%%%%%%%%%%%%%%%%%%%%%%%%%%%%%%%%%%%%%%%%%%%%%%%%%%%%%%
%% Chapter 4 — Architecture
%% Budget: 4 minutes / 3 slides
%%   S5  Architectural overview — high-level block diagram (Component / System)
%%   S6  Core abstractions      — Component, Port, System; unit-aware coupling
%%   S7  Master algorithms      — Gauss-Seidel / Jacobi / Hybrid / IJCSA
%%%%%%%%%%%%%%%%%%%%%%%%%%%%%%%%%%%%%%%%%%%%%%%%%%%%%%%%%%%%%%%%%%%%%%%%%%%%%

\section{Architecture}

%---------------------------------------------------------------------------
\begin{frame}{SysSimX Architecture}
    % TODO: overview diagram from thesis/figures/4_architecture/
    % One-line caption of what each layer does.
    \begin{center}
        \vspace{-1.5em}
        \includegraphics[height=\textheight,keepaspectratio]{4_architecture/syssimx_architecture.pdf}
    \end{center}
\end{frame}

%---------------------------------------------------------------------------
\begin{frame}{Core Abstractions: CoSimComponent}
    \begin{center}
        \vspace{-1.4em}
        \includegraphics[
            height=\textheight,
            keepaspectratio
        ]{4_architecture/cosimcomponent.pdf}
    \end{center}
\end{frame}

%---------------------------------------------------------------------------
\begin{frame}{Core Abstractions: System Assembly}
    \begin{center}
        \includegraphics[
            width=0.98\textwidth,
            height=0.92\textheight,
            keepaspectratio
        ]{4_architecture/system_architecture.pdf}
    \end{center}
\end{frame}

%---------------------------------------------------------------------------
\begin{frame}{Multi-Component Architecture}
    \centering
    \vspace{-0.35em}
    \begingroup
        \resizebox{!}{0.82\textheight}{%
            \begin{tikzpicture}[every node/.style={font=\sffamily}]
    % Public interface layer
    \begin{scope}[yshift=4.5cm]
        \node[layerbox, minimum width=10cm, minimum height=2cm] (publiclayer) at (4.5,0) {};
        \node[anchor=south west, font=\scriptsize\bfseries\itshape, text=gray!55!black]
            at ($(publiclayer.north west)+(0,-0.45)$) {public};

        \node[portbox, align=center] (In) at (0.75,0) {Input\\Ports};
        \node[portbox, align=center] (Out) at (8.25,0) {Output\\Ports};
        \node[portbox, align=center, minimum width=40mm] (interface) at (4.5, 0) {Unified\\Component Interface};

        \draw[arrow] (In) -- (interface);
        \draw[arrow] (interface) -- (Out);
    \end{scope}

    % Internal layer
    \begin{scope}
        \node[layerbox, fill=sxOrch!6, minimum width=10cm, minimum height=2cm] (internallayer) at (4.5,2) {};
        \node[anchor=south west, font=\scriptsize\bfseries\itshape, text=gray!55!black]
            at ($(internallayer.north west)+(0,-0.45)$) {internal};
        \node[inactivemodelbox, align=center] (modelA)   at (1, 2)  {Model A\\{\small(inactive)}};
        \node[activemodelbox, align=center]   (modelB)   at (8, 2)  {\textbf{Model B}\\{\small(active)}};
        \node[helperbox, align=center, line width=1.0pt]  (switching) at (4.5,2) {Mode\\Switching};
        \node[above, font=\sffamily\small, align=center] at (switching.north){selects active model};
        % --- Control relationships: switcher activates/deactivates models ---
        \draw[selectarrow] (switching.east) -- (modelB.west);
    \end{scope}

    % Switching support layer
    \begin{scope}[yshift=-0.5cm]
        \node[layerbox, fill=sxSwitch!6, minimum width=10cm, minimum height=2cm] (switchingsupport) at (4.5,0) {};
        \node[anchor=south west, font=\scriptsize\bfseries\itshape, text=gray!55!black]
            at ($(switchingsupport.north west)+(0,-0.45)$) {switching support};
        \node[helperbox, align=center] (selector)   at (2,   0) {Mode\\Selector};
        \node[helperbox, align=center] (hysteresis) at (4.5, 0) {Hysteresis};
        \node[helperbox, align=center] (adapter)    at (7,   0) {State\\Adapter};
    \end{scope}

    % --- Data path: interface delegates to whichever model is active ---
    \draw[arrow] (interface.south) --++(0, -0.75)
                                   --++(3.5, 0) node[midway, font=\sffamily\small, fill=white, inner sep=1pt] {delegate} 
                                   -- (modelB.north);
    \coordinate (supportJunction) at ($(switching.south)+(0,-0.75)$);

    \draw[line width=0.95pt, draw=sxSwitch!85!black]
    (switching.south) -- node[above left=2pt, font=\sffamily\small, inner sep=1pt] {uses} (supportJunction);

    \draw[supportarrow]
    (supportJunction) -- (hysteresis.north);

    \draw[supportarrow]
    (supportJunction) --++ (2.5, 0) -- (adapter.north);

    \draw[supportarrow]
    (supportJunction) --++ (-2.5, 0) -- (selector.north);

    % Helpers
    \useasboundingbox (-0.5, -1.5) rectangle (9.5, 5.5);
    % \draw[help lines, step=0.5, draw=gray!50, very thin] (-0.5,-1.5) grid (9.5, 5.5);


\end{tikzpicture}
%
        }
    \endgroup
\end{frame}

%---------------------------------------------------------------------------
\begin{frame}{Master Algorithms}
  \centering
  {\small\itshape Pluggable algorithm interface --- System and
   time-advancement strategy are separated.}
  \vspace{0.8em}

  \footnotesize
  \begin{tabular}{
      l
      >{\raggedright\arraybackslash}p{0.42\textwidth}
      >{\raggedright\arraybackslash}p{0.30\textwidth}}
    \toprule
    \textbf{Algorithm} & \textbf{Strategy} & \textbf{Use when} \\
    \midrule

    Jacobi &
    Prepare all interfaces $\to$ step all components. \newline
    \ac{SCC}-local solver for loops. &
    Loose coupling; \newline parallelism matters \\
    \addlinespace[3pt]

    Gauss--Seidel &
    Generation-wise: prepare $\to$ step per generation. \newline
    \ac{SCC}-local solver for loops. &
    \textbf{Default continuous}; \newline tighter coupling \\
    \addlinespace[3pt]

    \ac{IJCSA} &
    \textbf{Global} interface Newton solve --- one coupled residual
    for all zero-delay inputs $\to$ step all. &
    Many or strongly-coupled loops across the system \\
    \midrule

    Hybrid &
    Wraps a continuous algorithm; adds event detect / localize /
    dispatch. &
    \emph{Auto-selected} when any component declares an event indicator \\
    \bottomrule
  \end{tabular}

  \vspace{1em}
  \footnotesize
  \emph{Hybrid is automatically activated by the framework;
        the user only chooses the continuous algorithm.}
\end{frame}
